%% notes-tex/database/pputida-cell-state/sections/method.tex
%% [[experiments.database.pputida-cell-state]]
%% https://github.com/Mjvolk3/torchcell/tree/main/notes-tex/database/pputida-cell-state/sections/method.tex
%% SOURCE: hand-authored; the tables it describes are generated by
%% experiments/database/scripts/build_pputida_cell_state_table.py

\section{Where the numbers come from\secstatus{todo}}\label{sec:method}

Two layers, and the split is what makes the tables auditable. Every count is probed
from the served store. The row set and the last column of Table~\ref{tab:pcsrows} are
curation, held as data in the generating script so the judgment is versioned with the
numbers it sits beside.

\notesec{The stores read} The probe reads the dev build tree,
\file{$DATA_ROOT/data/torchcell/<dataset>/processed/}, which is the tree a loader
writes. It is not the knowledge-graph tree under \file{database/data/torchcell/}, which
the build user owns and which is read-only here. The two were checked against the build
manifest of release 2026.10.17-4295f7ca: every record count in
Table~\ref{tab:pcsrows} equals the manifest's \texttt{n\_experiments} for that dataset,
which is the evidence that the probed stores are the served ones.

\notesec{Records are pickled, not JSON} An LMDB value is a pickle, so a probe reads it
with \texttt{pickle.loads}. Reading one as UTF-8 fails on the pickle protocol byte,
which is how this was found.

\notesec{Constants are interned} An \texttt{Environment}, a medium and a compound are
each stored once in a sibling \file{processed/interned/} LMDB, and a record carries
\texttt{\{"\$ref": <sha256>, "name": <hint>\}} in their place. A probe that does not
splice these back reads every environment as an opaque reference. The script splices
recursively before touching any field.

\notesec{Distinct environments are counted by value} The interned table is the
tempting route to the count and it is the wrong one: it holds every interned constant,
environments and media and compounds alike, and nothing guarantees each entry is still
referenced. So the probe reads every record, splices its environment and hashes the
spliced object, which makes two references to the same condition one environment. That
costs a full scan of 2.7 million records and it is exact. The result is cached in
\file{experiments/database/results/pputida_cell_state_probe.json} so rebuilding the
document does not rescan.

\notesec{The genome is on the reference, not the experiment} Every row's assembly in
Table~\ref{tab:pcsrows} is read from \sym{reference.genome\_reference}. The
experiment-level \texttt{genome\_reference} is null in all fifteen datasets. That is by
design rather than a defect: the genome never gates a record, so a record does not
carry one.

All fifteen pin the same assembly, \texttt{GCA\_000007565.2}, while the strain on that
assembly varies: KT2440 for the screens, and the engineered backgrounds IY1449b,
IPL400 and CJ522 for the production and tolerance work. An engineered strain's
episomal cassettes are carried on the genotype as perturbations rather than folded
into the assembly, which is why one accession covers every row.

\notesec{What the probe cannot see} Whether a transposon pool was barcoded, and
whether an alcohol was fed as a carbon source or applied as a stressor, are properties
of the experiment that no stored field distinguishes. Both are carried as curation, the
first in the italic line under the genotype column and the second in prose in
Sec.~\ref{sec:candidates}.
